\begin{ex}
	Cho hình chóp $S.ABCD$ có $ABCD$ là hình bình hành tâm $O$. Gọi $M, N$ lần lượt nằm trên các cạnh $SA, SC$ sao cho $SM = 2MA$; $SN = 2NC$. Trong các mệnh đề sau, mệnh đề nào đúng, mệnh đề nào sai?
	\choiceTFt
	{\True $SO$ cắt $\text{mp}(BMN)$ tại điểm $I$ $(I = MN \cap SO)$}
	{\True Giao tuyến của $(SAD)$ và $(BMN)$ là $MK$ $(K = SD \cap BI)$}
	{$CD$ cắt $\text{mp}(BMN)$ tại điểm $F$ $(F = CD \cap BN)$}
	{Các điểm $E, B, F$ không thẳng hàng, với $E$ là giao điểm của $AD$ và $MK$}
	\loigiai{
		\begin{itemchoice}
			\itemch \textbf{Đúng.} Trong mặt phẳng $(SAC)$, hai đường thẳng $SO$ và $MN$ cắt nhau tại $I$. Vì $MN \subset (BMN)$ nên $I = SO \cap (BMN)$.
			\itemch \textbf{Đúng.} Trong mặt phẳng $(SBD)$, gọi $K = SD \cap BI$. Khi đó $K \in SD \subset (SAD)$ và $K \in BI \subset (BMN)$, mặt khác $M \in SA \subset (SAD)$ và $M \in (BMN)$ nên $(SAD) \cap (BMN) = MK$.
			\itemch \textbf{Sai.} Hai đường thẳng $CD$ và $BN$ không cùng nằm trong một mặt phẳng (chéo nhau) nên không cắt nhau. Giao điểm $F$ của $CD$ với $(BMN)$ phải là giao điểm của $CD$ với $KN$ trong mặt phẳng $(SCD)$.
			\itemch \textbf{Sai.} Ba điểm $E, B, F$ cùng thuộc hai mặt phẳng phân biệt $(ABCD)$ và $(BMN)$ nên chúng phải thẳng hàng (cùng nằm trên giao tuyến của hai mặt phẳng đó).
		\end{itemchoice}
	}
\end{ex}